\chapter{Sustainability}
\label{ch:sustainability}

\todo[inline]{From the \textit{Project Guide}: \\ It addresses both the way sustainability is taken into account in the design
and the way the product or system contributes to sustainability.}

\todo[inline]{Chapter introduction pending. Remember to connect it to the previous chapters.}


\section{Environmental Sustainability}
\label{sec:environmental_sustainability}

One of the main focuses of research for HUGO would be the study of high aspect ratio wings. 

-Benefits of possible research

The main environmental benefits of this project would be the research the platform would enable. One of the main focuses of research for HUGO would be the study of high aspect ratio wings. 


-Benefits of high aspect ratio

-Benefits of versatility. The aircraft was designed with versatility in mind. By enabling the changing of wing-platforms it is possible to broadens the usability of the aircraft. This means that a small modification on an existing airframe can be done instead of the construction on an entirely new aircraft. 


-Benefits of the reduction of downtime.
Another benefit of the versatility of the design is the potential to reduce the development time for fuel-saving technologies. A fast implementation will lead to a larger overall reduction in aviation emissions.  

It was chosen early on that the aircraft would be entirely powered by SAF. Using SAF would reduce the net-emissions for the operation of the aircraft. Even with the relative small scale of the aircraft comparing to commercial aircraft if was decided to try to minimize any impact. This was not only done in an effort to minimize any impact impact of the aircraft but also with the belief that with the experience and knowledge gained would contribute to the overall implementation of SAF. 

-Benefits of the long lifetime. 


\section{Social Sustainability}
\label{sec:social_sustainability}

One of the major incentives for the project was the possibility of the democratization of high level aerospace research. The phenomenon which can be investigated with this aircraft was traditionally only reserved for very large and expensive wind-tunnels or very complex purpose built prototypes.

As described in the certification, logistics and manufacturing section (see reference) it takes a lot of resources to build and operating a prototype aircraft. Most of these resources do not go to the implementation of novel technologies but instead the application of existing components and methods. This causes a lot of time and resources to be spend on non-research activities. In the early littérateur review of the project it became apparent it had mostly been large organisations such as NASA which had been able to build and operate prototypes for just investigating one specific phenomenon. HUGO will enable small aerospace institutions and by that widen the scope of influx of ideas for novel configurations and designs.  
 
With its speed, relative small size and low production cost the drone can also be modified to be used as a weapon. This was something that under no circumstance should happen. Therefore early one it was decided that the ownership of every produced drone should be kept to management of HUGO. It would not be possible for anyone to buy the drone itself. In that way it is ensured that the drone would only be used for scientific purposes.  

Throughout the potential noise generated by HUGO was discussed. It was decided that if the noise generated by the aircraft exceed 100 decibels for any bystanders, modifications would be implemented to mitigate the noise. Therefore the noise of the aircraft was modelled and calculated for the mission profile. 

-Emissions. Certain emissions generated by a turbine is unhealthy for humans. This would be emissions such as sod and carbon-monoxide. Part of the choice of the airport and most of the operation being over the ocean was driven by the intention to minimise the proximity to high population areas.    

Both of these aspects were investigated to ensure that the design would be well within the EU Environmental certification (see certification)   

-Safety. Flutter is a dangerous phenomenon during flight and has been responsible for several crashes throughout history\cite{nasa2011modeling}. Therefore it is relevant to understand the phenomena to ensure safe operation of passenger flights. The ability to accurately estimate the behaviour of wing-platforms will enable the design of safer aircraft.  

Another element is the integration of novel control systems such as gust alleviation. This does not only increase the design freedom of the aircraft but could also potentially dampen out turbulence which often is a nuisance for passengers.   


\begin{comment}
-The long lifetime of airframes. The long operation life of the 
    
\end{comment}

\section{Economic Sustainability}
\label{sec:economic sustainability}


This product remain unique positioned in the market. It has the potential to disrupt the traditional wind-tunnel test market. As no other companies can offer a comparable service. This enables a freedom to set the pricing of the services.   

A detailed market analysis was made for the system which can be found (insert reference) the overall conclusion was that the production and subsequent operation of HUGO would be able to generate a profit for the operators.  



\section{Life Cycle Assessment}
\label{sec:LCA}
           
At this point it possible to make the final LCA
With the design, logistics and operations of HUGO being finalised. Having this it is now possible to modify the preliminary LCA to make the final LCA. 

The life of the aircraft was split into three overall separate processes: the production of the aircraft, the operation of the aircraft and finally end-of-life of the aircraft. Each of those processes would consist of sub-processes. Every process takes one (or several) input and produces an output.

Instead of starting from scratch for every component and material use in the LCA the Swiss Federal Administration LCI Database (BAFU:2025) was used. This holds detailed descriptions and defined process for 11000 different components and processes. 

Describe the approaches for missing components 
Certain components and materials was not present in the BCA database. In some cases a similar component was chosen and in other information was sourced and a entirely new process was defined


Insert openLCA diagram: 

 For the production of the aircraft the inputs would be all the subsystems: electrical, wings, fuselage, empennage, turbine, landing gear and pylons. The combination of these lead to the output of an aircraft. This is then used as input for the operation of the aircraft. For the operation fuel, transportation, maintenance and electricity was also used as input. For the End of life operation it was the scraps from the aircraft which was inputs. 

Using the preliminary LCA it was possible to identify that operations-phase and more specifically the production and the combustion of the fuel. The first choice of fuel was a SAF HEFA based on soybean imported from Brazil. After performing a preliminary LCA it became apparent that the emissions connected to the production of soybeans would be larger than the emissions captured by the soybeans. This is due to the extensive deforestation in Brazil done to make room for soybean production\footnote{\href{https://www.sei.org/features/brazilian-soy-exports-and-deforestation/}{Brazil}(12.06.2026)}. Therefore it was decided to switch to rapeseed oil which can be locally production in Europe on existing fields. The same method for the refinement of the produce into SAF was utilised. 

Finally using this principle it was possible to create a LCA consisting of 4343 individual processes and with 52547 process links. 

Describe the model the Intergovernmental Panel on Climate Change (IPCC) 2021
(ISO 14067) 

(insert digram of emissions) 

See diagram of max emissions. 

\todo[inline]{TBD any other sections we deem necessary.}











\todo[inline]{Chapter conclusion pending. Remember to connect it to the chapters that come after.}